% ============================================================================
% AUTHOR NOTES -- not part of the compiled paper. Every open item lives here,
% once, instead of scattered \todo markers through the prose. Cross-reference
% is PLAN.md's milestones (M1-M5) in the skbuild repo.
%
% - Abstract + Conclusion: need the headline evaluation numbers (M1-M3
%   acceptance + M4 instrumentation) before these can be written for real.
% - RQ0-RQ4 (\Cref{sec:eval}): each needs its actual data (M4's before/after
%   measurement, M1-M3's per-project runs, M4's oracle-suite coverage census,
%   M4's profiling, M1-M3's predicted-vs-actual reconciliation).
% - \cite{nguyen2020kbuild}: read the actual paper before submission and
%   correct \Cref{sec:intro}/\Cref{sec:related}'s characterization of its
%   execution strategy and scale -- currently inferred, not verified.
% - \Cref{sec:related}: add kmax/variability-Kbuild literature (check
%   run/README.org's "Existing works" note, in git history) once identified.
% - \Cref{sec:design:traversal}: expand once PLAN.md M2's root-invocation-
%   context work (ARCH/CROSS_COMPILE, real root-Makefile reachability) lands.
% - Discussion (\Cref{sec:eval:threats}): update the shell/eval and Kconfig-
%   validity sentences once PLAN.md M2's P3/P5 work lands.
% - Artifact footnote/URL: fill in once a public repo or Zenodo archive
%   (Dynaplex's approach) is decided.
% ============================================================================

\documentclass[acmsmall,screen,review]{acmart}
\usepackage{algorithm}
\usepackage{algpseudocode}
\usepackage{booktabs}
\usepackage{xcolor}
\usepackage[capitalize]{cleveref}

\newcommand{\todo}[1]{\textcolor{red}{\textbf{[TODO: #1]}}}

\AtBeginDocument{%
  \providecommand\BibTeX{{Bib\TeX}}}

\setcopyright{acmlicensed}
\copyrightyear{2027}
\acmYear{2027}
\acmDOI{10.1145/nnnnnnn.nnnnnnn}
\acmJournal{PACMSE}
\acmVolume{1}
\acmNumber{FSE}
\acmArticle{1}
\acmMonth{7}

\begin{document}

\title{skbuild: Extracting Kbuild Build Conditions via Eager-Merge Symbolic Execution}

\author{Anonymous Author(s)}
\affiliation{%
  \institution{Anonymous Institution}
  \country{USA}
}

\begin{abstract}
Large configurable systems such as the Linux kernel, BusyBox, and coreboot
describe what gets compiled with Kbuild-style Makefiles: a per-directory
\texttt{obj-\$(CONFIG\_X)}-shaped vocabulary layered on GNU Make, driven by a
Kconfig-defined space of Boolean and tristate options. A single build under
one \texttt{.config} answers only which files were selected for \emph{that}
configuration; it does not say how the build varies as options change, which
is what a developer needs to explain a crash under an unfamiliar
configuration, to bound the effect of flipping one option, or to identify
files a Kbuild refactor silently stopped building. We present skbuild, a
symbolic-execution engine for the Kbuild-relevant subset of GNU Make that
computes, for every object file, a Z3-checkable Boolean condition over
Kconfig symbols describing exactly when that file participates in the
build. Following a fork-per-branch design that forks a new execution path
at every \texttt{ifeq}/\texttt{ifdef} makes real Kbuild trees intractable ---
a 30-line synthetic fixture already produces dozens of forked paths, and a
first attempt at the checked-in BusyBox snapshot exhausts the pattern
entirely before finishing --- skbuild instead keeps a single symbolic state
per Kbuild file, representing each Make variable as a set of
\texttt{(word, condition)} pairs, and merges a conditional's two branches
back into that one state immediately, gating each branch's contribution by
its own guard. This design is directly adapted from Cybolic's eager-merge
representation for CMake~\cite{cybolic}, applied here to a build-description
language with a different, and in some ways simpler, variability model:
Kbuild's symbolic inputs are Kconfig's small, explicitly finite Boolean and
tristate domains rather than CMake's unbounded-string \texttt{option()}
values. We evaluate skbuild on \todo{BusyBox / coreboot / Linux, N pinned
releases, M reconciled object sets, K categorized mismatches}.
\end{abstract}

\begin{CCSXML}
<ccs2012>
<concept>
<concept_id>10011007.10011074.10011099.10011692</concept_id>
<concept_desc>Software and its engineering~Software verification</concept_desc>
<concept_significance>500</concept_significance>
</concept>
<concept>
<concept_id>10011007.10011074.10011099</concept_id>
<concept_desc>Software and its engineering~Software testing and debugging</concept_desc>
<concept_significance>300</concept_significance>
</concept>
</ccs2012>
\end{CCSXML}

\ccsdesc[500]{Software and its engineering~Software verification}
\ccsdesc[300]{Software and its engineering~Software testing and debugging}

\keywords{build systems, Kbuild, Linux kernel, symbolic execution, build conditions, SMT solving, variability}

\maketitle

\section{Introduction}
\label{sec:intro}

Large, configurable C codebases --- the Linux kernel, BusyBox, coreboot, and
other projects that adopted the Kbuild convention --- do not compile a fixed
set of source files. A per-directory \texttt{Makefile} or \texttt{Kbuild}
file assigns object names to variables such as \texttt{obj-y}, \texttt{obj-m},
and \texttt{lib-y} conditionally on Kconfig options
(\texttt{obj-\$(CONFIG\_USB) += usb.o}), nests those assignments inside
\texttt{ifeq}/\texttt{ifdef} blocks, and recurses into subdirectories whose
membership is itself conditional. Kconfig, in turn, defines each
\texttt{CONFIG\_*} symbol's domain: most commonly Boolean
(\texttt{y}/not-set) or tristate (\texttt{y}/\texttt{m}/not-set), sometimes a
small enumerated or numeric range. A single invocation of \texttt{make} with
one \texttt{.config} answers a narrow question --- which files were selected
for \emph{that} configuration --- but does not directly answer how the build
varies across the configuration space the project actually ships.

That broader question matters in practice. A crash report from a shipped
kernel build names a source file, but that file may or may not have been
compiled in depending on options the reporter never mentioned; knowing
\emph{which} configurations could have produced the crashing binary narrows
root-cause search before any other analysis starts. A patch that touches a
file guarded behind \texttt{CONFIG\_DEBUG\_INFO} needs a build configuration
under which the file is even reachable before static analysis of the patch
can begin. A maintainer refactoring a Kbuild tree wants to know whether the
refactor changed which files are selected under any configuration, not just
under whichever one happens to be `make defconfig`. All three reduce to the
same question this paper answers: \emph{given an object file, under what
Boolean condition over Kconfig symbols is it selected by the build?}

Nguyen and Nguyen~\cite{nguyen2020kbuild} showed that this question is
answerable by symbolic execution of Kbuild Makefiles specifically, treating
\texttt{CONFIG\_*} variables as finite-domain symbolic inputs and Z3 as the
underlying solver. This paper's skbuild is a from-scratch, from-first-principles
reimplementation of that direction: it keeps the CONFIG-symbols-as-symbolic-
inputs framing and the Z3 backend, but replaces a fork-per-branch execution
strategy --- which we found, empirically, explodes even on small synthetic
fixtures and cannot complete on a real multi-directory tree at all --- with
an eager branch-and-merge representation adapted from Cybolic, a recent
system that solved the structurally identical problem (extract a per-file
build condition via symbolic execution of a configuration-driven build
description) for CMake~\cite{cybolic}. Kbuild's variability model is, in one
concrete way, simpler than CMake's: Kconfig options are explicitly
finite-domain by construction (Boolean or tristate, occasionally a small
enumerated set), where CMake's \texttt{option()}/cache variables are
unbounded strings that Cybolic must itself constrain to a finite domain
before Z3 can reason about them. Kbuild's variability model is, in another
way, its own challenge: recursive directory traversal is itself
conditional, so the same eager-merge idea has to apply not just to
variable values within one file but to \emph{which files get analyzed at
all}.

\emph{Contributions.} In summary, this paper makes the following
contributions:

\begin{itemize}
\item We identify, with a concrete measurement on a small synthetic
  fixture and a real-world corpus, that fork-per-branch symbolic execution
  of Kbuild Makefiles is not viable at real-project scale (\Cref{sec:overview}).
\item We adapt Cybolic's eager branch-and-merge symbolic-state
  representation~\cite{cybolic} to Kbuild/GNU Make, including two
  Kbuild-specific correctness subtleties --- conditionally-named
  \texttt{:=} overwrites, and re-gating carried-forward (untouched) values
  by branch guard at merge time --- that do not arise in CMake's simpler
  assignment model (\Cref{sec:design}).
\item We implement this design in skbuild, a pure-Python, Z3-backed
  analyzer with no other runtime dependency (\Cref{sec:eval}).
\item We evaluate skbuild on BusyBox, coreboot, and the Linux kernel,
  reconciling predicted object sets against real, independently obtained
  build output and categorizing every mismatch (\Cref{sec:eval}).
\end{itemize}

Tool and evaluation artifacts \todo{artifact URL}.

\section{Overview}
\label{sec:overview}

\subsection{A Worked Example}
\label{sec:overview:example}

\Cref{fig:example} shows a small Makefile exercising the constructs this
paper is about: Boolean-guarded object selection
(\texttt{obj-\$(CONFIG\_1)}), an unconditional overwrite whose \emph{target
variable name} is itself conditional (\texttt{obj-\$(CONFIG\_A) := 1.o}),
and a variable (\texttt{BITS}) reassigned unconditionally by both arms of an
\texttt{ifeq}/\texttt{else}.

\begin{figure}
\begin{verbatim}
obj-$(CONFIG_1) += 1.o
obj-$(CONFIG_A) := 1.o        # overwrites obj-y's value *when CONFIG_A=y*
obj-$(CONFIG_A) += 3.o
obj-$(CONFIG_B) += 3.o

ifeq ($(CONFIG_A),y)
  BITS := 32
else
  BITS := 64
endif
obj-y += probe$(BITS).o
\end{verbatim}
\caption{A worked example. \texttt{obj-y}'s final value for \texttt{1.o}
is \texttt{CONFIG\_A=y OR (CONFIG\_1=y AND CONFIG\_A!=y)}: the first
\texttt{+=} contributes \texttt{1.o} whenever \texttt{CONFIG\_1=y}, and the
later \texttt{:=} only erases that contribution in the universe where it
actually fires, i.e.\ \texttt{CONFIG\_A=y}. \texttt{BITS} ends with exactly
two values, \texttt{32} under \texttt{CONFIG\_A=y} and \texttt{64}
otherwise --- \texttt{0} does not survive as a spurious third case, even
though it would if the merge at the \texttt{ifeq} naively unioned each
branch's raw carried-forward state.}
\label{fig:example}
\end{figure}

A fork-per-branch symbolic executor treats every \texttt{ifeq} as a point
where execution splits into two independent copies of the current state,
one per branch, carried forward until (if ever) they are recombined. Under
this design, the seven statements in \Cref{fig:example} already fork
32 distinct execution paths in our reimplementation of that strategy
(\Cref{sec:eval:motivation}), and the checked-in BusyBox snapshot's
top-level tree does not finish analysis at all under it.

skbuild instead keeps exactly one symbolic state per Kbuild file throughout.
Each Make variable's value is a map from word to the Z3 condition under
which that word is currently a member (\Cref{sec:design:rep}). A conditional
clones this one state per branch, runs each branch's statements against its
own clone, and immediately folds both clones back into the original state
before the next statement (\Cref{sec:design:merge}) --- so state size stays
proportional to program size, not to the number of branches taken so far.

\subsection{Design Summary}
\label{sec:overview:design}

skbuild has four pieces: a GNU Make parser (a Python 3 port of Mozilla's
\texttt{pymake}), a symbolic evaluator over the guarded-value state
described above, a Z3-backed condition solver, and a directory-traversal
layer that recurses into subdirectories under their own guard, using the
same eager-merge state to decide which subdirectories to visit at all.
Recipes and rules are parsed but never executed; \texttt{shell} is
disabled by default (\Cref{sec:eval:threats} discusses this limitation).
The tool reports one record per selected object: a path, a
target kind (built-in / module / library), and a Boolean condition over
Kconfig symbols.

\section{Design}
\label{sec:design}

\subsection{Representing Branch-Dependent Values}
\label{sec:design:rep}

A Make variable in skbuild (\texttt{ds.VarG}) is a name, a flavor
(\texttt{RECURSE} for \texttt{=}/\texttt{define}, \texttt{SIMPLY} for
\texttt{:=}/\texttt{::=}/\texttt{+=}), and a map from word to Z3 condition.
A symbolic state (\texttt{ds.BaseState}) is a map from variable name to
\texttt{VarG}. This is the direct Kbuild analogue of Cybolic's
\texttt{Switch}/\texttt{Const}/\texttt{Concat} value representation for
CMake strings and lists~\cite{cybolic}, specialized to the fact that a
Kbuild variable's value is, in every construct this paper's supported
subset covers, a whitespace-separated set of independently-conditioned
words rather than an order-sensitive string.

An assignment updates this state under an ambient guard --- the conjunction
of every enclosing conditional's branch condition --- threaded as an
explicit parameter through statement execution rather than stored per-path.
A \texttt{+=} accumulates: each newly assigned word's condition is OR'd
with the guard into its existing entry, or created fresh. A \texttt{:=} or
\texttt{=} overwrites, but --- and this is the first of two correctness
subtleties that do not arise in CMake's simpler \texttt{set()} --- only
\emph{within} the guard under which this particular assignment's name and
value actually resolved. Kbuild's \texttt{obj-\$(CONFIG\_A) := 1.o} example
in \Cref{fig:example} resolves to the variable literally named
\texttt{obj-y} only when \texttt{CONFIG\_A=y}; outside that condition, this
statement assigns to \texttt{obj-undef} (a name nothing reports) and must
leave \texttt{obj-y}'s prior value untouched. skbuild's overwrite therefore
computes, per pre-existing word $w$ with condition $c$: $c \wedge
\lnot\mathit{cond}$ carried forward (the assignment did not fire), unioned
with $\mathit{cond}$ for each newly assigned word, where $\mathit{cond}$ is
this specific assignment's fully-resolved condition (guard $\wedge$ name
match $\wedge$ value match).

\subsection{Branch Execution and Eager State Merging}
\label{sec:design:merge}

A \texttt{ConditionBlock} clones the current state twice, executes the
\texttt{then} branch's statements against one clone under
\texttt{guard} $\wedge$ \texttt{if\_cond}, the \texttt{else} branch's
statements against the other under \texttt{guard} $\wedge \lnot$
\texttt{if\_cond}, and merges. The second correctness subtlety --- again,
one with no CMake analogue, because CMake's \texttt{if()}/\texttt{else()}
does not have Kbuild's pattern of an unconditional-within-the-branch
overwrite whose \emph{reachability} is what's conditional --- is that
merging cannot simply union each branch's raw per-word conditions computed
by \Cref{sec:design:rep}'s overwrite rule. Consider \Cref{fig:example}'s
\texttt{BITS}: entering the \texttt{ifeq}, \texttt{BITS} has no prior value;
each branch's \texttt{:=} computes a fresh, branch-guard-scoped condition
for its own word (\texttt{32} under \texttt{CONFIG\_A=y}, \texttt{64} under
its negation) --- unremarkable. But if \texttt{BITS} had already been set to
\texttt{0} before this \texttt{ifeq} (as it is, from an earlier
\texttt{ifeq (\$(CONFIG\_B),y) BITS := 0 endif}, in the fuller version of
this fixture we use for testing), each branch's own overwrite rule computes
a \emph{carry-forward} term for \texttt{0} of ``\texttt{CONFIG\_B=y} AND
this branch's assignment did not fire'' --- and since the two branches'
``did not fire'' conditions are exact complements of each other, naively
OR-ing the two branches' raw \texttt{0} terms together reconstructs
\texttt{CONFIG\_B=y} unconditionally, resurrecting a value that both
branches, between them, always overwrite. The fix is to re-gate each
branch's \emph{entire} per-word contribution by that branch's own guard at
merge time, not just rely on the per-statement conditions computed inside
the branch:
\[
\mathit{merged}(w) \;=\; \bigl(\mathit{guard}_{\mathit{then}} \wedge
\mathit{then\_state}(w)\bigr) \;\vee\; \bigl(\mathit{guard}_{\mathit{else}}
\wedge \mathit{else\_state}(w)\bigr)
\]
which is exactly what Cybolic's own \texttt{merge\_branches} does for CMake
values~\cite{cybolic}; we did not initially replicate this and had to find
the bug by hand-deriving \Cref{fig:example}'s expected conditions and
comparing against skbuild's output (\Cref{sec:eval:motivation} discusses
this as a methodology point, not just a bug-fix anecdote). Applying this
re-gating to \emph{every} variable at every conditional, however, grows
condition expressions unboundedly across a file regardless of relevance:
an untouched variable's condition would be re-wrapped in
$(\mathit{guard}_{\mathit{then}} \wedge \cdot) \vee
(\mathit{guard}_{\mathit{else}} \wedge \cdot)$ at every subsequent
conditional even though that expression is always tautologically equal to
the untouched value. skbuild tracks, per state, the set of variable names
actually reassigned since the state was cloned (\texttt{State.touched});
merge only rebuilds names in the union of both branches' touched sets, and
leaves everything else in the destination state exactly as it was ---
mirroring Cybolic's own \texttt{changed\_vars} restriction to the same
end~\cite{cybolic}.

\subsection{Recursive Traversal}
\label{sec:design:traversal}

Kbuild's \texttt{obj-\$(CONFIG)} pattern applies to subdirectory names too
(\texttt{obj-\$(CONFIG\_USB) += usb/}); skbuild's \texttt{BaseState}
aggregates, across every non-target, non-ignored variable in the final
state, a map from subdirectory to the Z3 condition under which it is
reached, OR-ing conditions where more than one variable names the same
subdirectory. Today this traversal is a straightforward per-directory
recursion; it does not yet model \texttt{ARCH}/\texttt{CROSS\_COMPILE} or
derive root-directory reachability from the real root Makefile rather than
a curated top-directory list (\Cref{sec:eval:threats}).

\section{Evaluation}
\label{sec:eval}

\subsection{RQ0 (Methodology): Fork-Per-Branch Does Not Scale}
\label{sec:eval:motivation}

\todo{This subsection reports the actual measurement instead of the
anecdotal numbers used in \Cref{sec:overview}: re-run the retired
fork-per-branch implementation (available in git history before this
paper's rewrite) against \Cref{fig:example} and the BusyBox/Linux
snapshots, and report path counts / wall time / where it fails outright,
as a controlled before/after rather than a one-off observation.}

\subsection{RQ1: Applicability}
\label{sec:eval:rq1}

\todo{Does skbuild run to completion on a pinned BusyBox release, a pinned
coreboot release, and a pinned Linux kernel release, per PLAN.md's M1--M3
acceptance criteria? Report per-project: source LOC, number of
Makefile/Kbuild files, wall time, and completion status.}

\subsection{RQ2: Coverage}
\label{sec:eval:rq2}

\todo{Per-construct usage census across BusyBox/coreboot/Linux (PLAN.md's
M4), classified the way Cybolic's RQ2 classifies CMake commands: (i)
\emph{modeled}, (ii) \emph{correctly out of scope} (affects recipes,
linking, or installation, not which files are selected --- a no-op is
provably harmless), or (iii) \emph{deliberately unsupported}, reported
with real usage-frequency counts rather than a bare "N warnings" number,
so a reader can see whether an unmodeled construct is a corner case or, as
with Cybolic's \texttt{get\_filename\_component} finding, the single most
frequently invoked construct across every project evaluated. \texttt{shell}
probes of the host/toolchain are the leading candidate for that role in
Kbuild; state plainly whether that turned out true.}

\subsection{RQ3: Performance}
\label{sec:eval:rq3}

\todo{Wall time, peak RSS, and Z3 call counts across the three projects,
per PLAN.md's M4 instrumentation. State explicitly whether skbuild's
dominant cost is Z3 solver calls under deep conditional nesting or
something else --- do not assume it mirrors Cybolic's string/list
concatenation bottleneck without measuring.}

\subsection{RQ4: Agreement with Real Builds}
\label{sec:eval:rq4}

This is the paper's central empirical result. For each project, we use Z3
to extract, from skbuild's own reported conditions, a small set of concrete
witness configurations that between them are predicted to cover every
reported object (a greedy set cover over conditions, the same technique
underlying Cybolic's sufficient-CI-matrix result~\cite{cybolic}), plus at
least one negative witness per file where feasible. We materialize each
witness as a real \texttt{.config}, run the project's actual build against
it in a clean checkout, and compare the real build's object/module output
against skbuild's prediction for that same witness. \todo{Report, per
project: number of witnesses, and true/false positive/negative object
counts, with every non-match causally categorized (missing generated file,
unsupported construct, genuine skbuild bug, or a witness Kconfig itself
would reject) per PLAN.md's M1.0 tool --- not a bare diff count.}

\subsection{Discussion and Threats to Validity}
\label{sec:eval:threats}

skbuild does not execute \texttt{shell} or \texttt{eval} side effects by
default (\Cref{sec:overview:design}); any file whose selection genuinely
depends on such an effect is reported as an incomplete result, not silently
omitted or silently assumed absent. This is a real completeness gap, not
just a conservative default: a future opt-in effect boundary (a sandboxed
command runner) is planned but not yet implemented, and every RQ2/RQ4
number in this paper should be read against that gap. Kconfig validity ---
distinguishing a condition that is merely Make-satisfiable from one a real
\texttt{.config} could actually produce under Kconfig's own
dependency/select/choice rules --- is not modeled at all today: skbuild
reports the former, not the latter, so a reported condition can describe a
configuration Kconfig itself would reject. Every evaluation number in this
paper is tied to one pinned source
revision per project; Kbuild conventions have historically drifted across
kernel releases (see the retired \texttt{LINUX\_ANALYSIS\_TODO.md}'s own
documented gaps between kernel versions in this project's history), so
these numbers should not be read as claims about every release of these
projects.

\section{Related Work}
\label{sec:related}

\emph{Symbolic execution of build descriptions.} Nguyen and
Nguyen~\cite{nguyen2020kbuild} first applied symbolic execution
specifically to Linux Kbuild Makefiles, treating \texttt{CONFIG\_*}
variables as finite-domain symbolic inputs solved with Z3; this paper is a
from-scratch reimplementation of that direction under an eager
branch-and-merge execution strategy (\Cref{sec:design}) rather than that
system's own strategy.
Tamrawi et al.'s SYMake~\cite{tamrawi2012symake} symbolically evaluates GNU
Make itself (not specifically Kbuild) to analyze build code more generally.
King's original symbolic execution~\cite{king1976symbolic} underlies both
lines of work, applied here to a declarative build description rather than
an imperative program.

\emph{CMake.} Cybolic~\cite{cybolic} solves the structurally identical
problem --- per-file build condition extraction via symbolic execution ---
for CMake, and is the direct source of this paper's eager branch-and-merge
representation (\Cref{sec:design}); \Cref{sec:intro} discusses where
Kbuild's variability model is simpler (finite Kconfig domains) and where it
is not (conditional recursive traversal). Nguyen, Nguyen, and
Phan~\cite{nguyen2022analyzing} posed the CMake build-condition question
that Cybolic answers, as a vision piece with neither implementation nor
evaluation.

\emph{Variability-aware analysis.} Erwig and Walkingshaw's choice
calculus~\cite{erwig2011choice} provides a general representation for
software variation that this paper's guarded-value state can be understood
as a build-system-specific instance of.

\section{Conclusion}
\label{sec:conclusion}

We presented skbuild, a symbolic-execution engine that computes, for every
object file in a Kbuild-style build, a Z3-checkable Boolean condition over
Kconfig symbols describing when that file is selected. Its central design
decision --- an eager branch-and-merge symbolic state, adapted from
Cybolic's CMake analyzer, in place of a fork-per-branch design we found
does not scale to real Kbuild trees --- required two Kbuild-specific
correctness fixes with no CMake analogue: scoping a conditionally-named
overwrite to the universe in which it actually fires, and re-gating
carried-forward values by branch guard at merge time.
\todo{headline evaluation result}

\bibliographystyle{ACM-Reference-Format}
\begin{thebibliography}{99}

\bibitem{cybolic}
Anonymous Author(s). 2027.
Cybolic: Extracting CMake Build Conditions via Symbolic Execution.
\emph{Proc. ACM Softw. Eng.}, FSE (2027).

\bibitem{nguyen2022analyzing}
ThanhVu Nguyen, KimHao Nguyen, and Anh Cuong Phan. 2022.
Analyzing the CMake Build System.
In \emph{Proceedings of the ACM/IEEE 44th International Conference on Software Engineering: Software Engineering in Practice (ICSE-SEIP '22)}.

\bibitem{nguyen2020kbuild}
KimHao Nguyen and ThanhVu Nguyen. 2020.
Using Symbolic Execution to Analyze Linux KBuild Makefiles.
In \emph{Proceedings of the IEEE International Conference on Software Maintenance and Evolution (ICSME '20)}.

\bibitem{tamrawi2012symake}
Ahmed Tamrawi, Hoan Anh Nguyen, Hung Viet Nguyen, and Tien N. Nguyen. 2012.
Build Code Analysis with Symbolic Evaluation.
In \emph{Proceedings of the 34th International Conference on Software Engineering (ICSE '12)}.

\bibitem{king1976symbolic}
James C. King. 1976.
Symbolic Execution and Program Testing.
\emph{Commun. ACM} 19, 7 (1976).

\bibitem{erwig2011choice}
Martin Erwig and Eric Walkingshaw. 2011.
The Choice Calculus: A Representation for Software Variation.
\emph{ACM Trans. Softw. Eng. Methodol.} 21, 1 (2011).

\end{thebibliography}

\end{document}
